\chapter{JSFX scripts}
\label{ch:jsfx}

\fxme\ can run JSFX scripts, the effects written in REAPER's EEL2 language,
in the server's chain, next to VST3 plugins. The scripts live on the computer
that runs the DAW: when you add one, the plugin packs it with what it needs
and sends it to the server, which runs it with its own copy of the JSFX
engine. Nothing has to be installed on the server, and REAPER is not needed
on either computer.

\section{Where scripts come from}

The plugin looks for scripts in REAPER's \file{Effects} folder on this
computer, and in the folders you add. The \ui{JSFX SCRIPTS} part of the
\ui{Settings} tab shows how many it found and where it searched.
\ui{Folders\ldots} adds or removes a folder; \ui{Rescan} looks again after you
have installed or changed scripts.

An added folder is laid out as REAPER's: scripts in subfolders of an
\file{Effects} folder, their data files (samples, impulse responses, amp
models) in a \file{Data} folder next to it.

\section{Adding a script}

Scripts appear in \ui{SERVER PLUGINS} after the server's plugins, tagged
\ui{JS}. The filter also matches their description. Add one as you would a
plugin: it then behaves as any slot of the chain (order, bypass, routing,
latency compensation, automation, saved with the project).

When you add a script, the plugin sends with it:

\begin{itemize}
  \item the script and the files it imports;
  \item the data files of its file sliders and the files it names in its
    code, such as images;
  \item audio data files the engine cannot read itself (FLAC, AIFF, Ogg,
    MP3, 32-bit or floating-point WAV), converted to 24-bit WAV on the way.
\end{itemize}

Data and images may take up to 48~MB, the whole script 56~MB. Everything
sent is saved in the DAW project with each instance of the script, so the
project opens again even if the script has changed or gone from this
computer. To take a newer version of a script, remove it from the chain and
add it again.

\section{Controls}

A script's sliders are its parameters. In the generic controls they show in
the script's own units; a slider with named values, or one that chooses a
file, is a list. Sliders whose name starts with a minus sign are hidden, as
in REAPER: they can still be automated and mapped, but have no row.

A script that reports a delay (\texttt{pdc\_delay}, in lookahead effects) has
it compensated like a plugin's latency.

\section{The script's display}

A script with its own display (\texttt{@gfx}) shows it in the
\ui{Parameters} tab, under its sliders, as REAPER's effect window does.
\ui{Generic controls} switches to the sliders alone.

The display is drawn on this computer from a copy of the server's script,
kept in step with it: what the audio code computes (spectra, meters, curves)
comes from the server, at the telemetry rate. The mouse, the wheel and the
keys work as in REAPER, and so do the script's menus and mouse cursors.
What you change on the display, a slider or the script's own data, goes to
the server.

You can drop files on a display that accepts them. The script receives the
path on \emph{this} computer, so a script that opens the dropped file on the
server will not find it.

\section{Limits}

\begin{itemize}
  \item No presets (\file{.rpl} files) and no script editor: edit scripts
    in REAPER or a text editor, then remove and add the script again.
  \item The scripts of all sessions on a server share one \texttt{gmem[]},
    as the scripts of one REAPER project do.
  \item A script's display may differ from REAPER's in small ways: it may
    keep redrawing while no audio flows, and a menu that a script opens by
    itself (without a click or a key) cannot be answered.
  \item Some display variables may lag when the script's audio code sets
    them only from within functions. Tell us which script if you see one
    whose display stays still.
\end{itemize}
